%% ma: L12-B04-0017
%% lop: 12
%% chuong: 1
%% bai: 4
%% dang: 12.4.14
%% loai: DS
%% muc: [NB, TH, TH, TH]
%% dap_an: "ĐSĐS"
%% nguon: "HK1-DeOn_1_1.pdf (M10, Đề ôn tập học kì 1 Toán 12, Đề số 2), Phần II Câu 1"
%% kien_thuc: "Bài 1 (tính đơn điệu và cực trị), Bài 3 (đường tiệm cận), Bài 4 (khảo sát và vẽ đồ thị hàm số)"
%% trang_thai: chinh-thuc
%% ghi_chu: "Ý c gốc ghi \"Trục đối xứng của đồ thị hàm số đi qua điểm (−4; −4)\"; SGK KNTT chỉ nêu tâm đối xứng của đồ thị hàm phân thức. Bản kho sửa thành \"Tâm đối xứng ... là điểm (−4; −4)\", đáp án Đúng (khớp ý đồ gốc). Đã duyệt 29/10/2026."
\begin{ex}
Cho hàm số $y=f(x)=\dfrac{x^2+4x+1}{x+4}$.
\choiceTF{\True Tập xác định của hàm số đã cho là $D=\mathbb R\setminus\{-4\}$}{Hàm số đã cho không có cực trị}{\True Tâm đối xứng của đồ thị hàm số là điểm $(-4;-4)$}{Đồ thị hàm số có một đường tiệm cận xiên và một đường tiệm cận ngang}
\loigiai{a) Đúng.

b) Sai: $f'(x)=\dfrac{x^2+8x+15}{(x+4)^2}=\dfrac{(x+3)(x+5)}{(x+4)^2}$ đổi dấu tại $x=-5$, $x=-3$, hàm số có hai điểm cực trị.

c) Đúng: $f(x)=x+\dfrac1{x+4}$ có tiệm cận đứng $x=-4$, tiệm cận xiên $y=x$; tâm đối xứng là giao điểm $(-4;-4)$.

d) Sai: đồ thị có tiệm cận đứng $x=-4$ và tiệm cận xiên $y=x$, không có tiệm cận ngang.}
\end{ex}
